% ================================================================
% A Comparative Study of 1D and 2D Convolutional Architectures
% for Bird Sound Classification under a Microcontroller Flash Budget
% IEEE Conference Paper Format
% All numbers from verified_results.csv (canonical configs only)
% ================================================================

\documentclass[conference]{IEEEtran}
\usepackage[utf8]{inputenc}
\usepackage{cite}
\usepackage{amsmath}
\usepackage{amssymb}
\usepackage{graphicx}
\usepackage{booktabs}
\usepackage{hyperref}

\hyphenation{micro-controller}

\begin{document}

\title{A Comparative Study of 1D and 2D Convolutional Architectures for Bird Sound Classification under a Microcontroller Flash Budget}

\author{Muhammad Mun'im Ahmad Zabidi%
\thanks{Southeast Asian Institution}}

\maketitle

\begin{abstract}
We systematically compare temporal convolutional networks (TCNs) and 2D CNNs for 12-class bird sound classification on an Arduino Portenta H7 (Cortex-M7, 512\,KB flash budget). Evaluating 8 TCN ablation variants (3 seeds each, INT8), a raw-waveform 1D CNN, a 2D CNN baseline, and EfficientNetB0 transfer learning on the source-disjoint MyGardenBird dataset, we find: (i)~the best in-budget TCN reaches \textbf{92.92\%} INT8 at 266\,KB and ties the 2D CNN baseline---but the 2D CNN requires 1{,}630\,KB (3$\times$ over budget) to match this accuracy; (ii)~among TCN variants that fit the 512\,KB budget, accuracy ranges 81--93\%, with competitive performance requiring residual/squeeze-excite depth; (iii)~EfficientNetB0 suffers 3.84\,pp INT8 loss due to unbounded swish activations; (iv)~ShuffleNetV2-Compact at 476\,KB achieves 88.66\% INT8, the only in-budget non-TCN alternative but 4.26\,pp below best TCN. Under strict flash constraints, 1D TCNs make better use of the deployable parameter budget than 2D approaches, contradicting the intuition that vision-based spectrograms always benefit from 2D convolutions.
\end{abstract}

\begin{IEEEkeywords}
temporal convolution, audio classification, microcontroller, TinyML, INT8 quantization, bird sound, budget-aware neural architecture
\end{IEEEkeywords}

\IEEEpeerreviewmaketitle

% ================================================================
% I. INTRODUCTION
% ================================================================

\section{Introduction}

Acoustic monitoring on battery-powered edge devices requires models that fit within strict microcontroller (MCU) constraints. Our target platform is the \textbf{Arduino Portenta H7} (ARM Cortex-M7, 480\,MHz), which provides a realistic deployable budget of \textbf{512\,KB flash} for model weights and runs \textbf{TensorFlow Lite Micro} with INT8 quantization. This is a hard constraint: a model exceeding 512\,KB cannot be deployed, regardless of accuracy. All architectural choices are therefore evaluated not only on accuracy but on whether they fit this budget.

The question we investigate is fundamental: \textit{Which architectural family---temporal convolutions (1D) or spatial convolutions on mel-spectrograms (2D)---makes more efficient use of a fixed flash budget for bird sound classification?}

The conventional wisdom favors 2D convolutions on spectrograms, treating the mel-spectrogram as a ``2D image'' and leveraging the spatial structure of time--frequency representations. However, this intuition was developed in regimes with no flash budget; under a strict 512\,KB constraint where parameter count directly determines deployability, the trade-off changes. TCNs operate on 1D sequences and may require less per-unit-accuracy overhead.

We address this question through a systematic comparison of four architectural families:
\begin{enumerate}
  \item Raw-waveform 1D CNN (no spectrogram frontend)
  \item Temporal convolutional networks (8 variants, 3 seeds, INT8)
  \item 2D CNN baseline on mel-spectrograms
  \item Transfer-learned EfficientNetB0 (ImageNet pre-trained)
\end{enumerate}

Additionally, we evaluate two in-budget peer architectures (ShuffleNetV2-Compact and SqueezeNet) to contextualize TCN performance on a Pareto frontier.

\textbf{Contributions:}
\begin{enumerate}
  \item A budget-aware comparison showing that the best in-budget TCN (92.92\% INT8, 266\,KB) ties the 2D CNN baseline in accuracy---but the 2D CNN requires 1{,}630\,KB (3$\times$ over budget) to achieve this, while the TCN fits the H7 flash limit.
  \item An 8-variant TCN ablation (3 seeds, INT8) showing that competitive $\sim$93\% accuracy across all families requires depth and residual connections.
  \item Evidence that, under tight flash budgets, 1D TCNs are more size-efficient than 2D CNNs for audio classification, contradicting prior assumptions.
\end{enumerate}

% ================================================================
% II. DATASET AND EXPERIMENTAL PROTOCOL
% ================================================================

\section{Dataset and Experimental Protocol}

\subsection{Target Platform and Budget Justification}

The Arduino Portenta H7 was selected for four reasons: (1)~its Cortex-M7 core (480\,MHz, with FPU and CMSIS-NN SIMD support) is representative of the upper end of the MCU class, so results transfer to a range of M7-based platforms; (2)~TensorFlow Lite Micro provides official, reproducible INT8 support and deployment tooling; (3)~its low-power operation (15\,mA per inference cycle) suits battery/solar remote sensors; (4)~while the H7 has 2\,MB total flash, realistic deployments reserve flash for bootloader, audio buffers, the mel-spectrogram front-end, and TFLM runtime, leaving \textbf{512\,KB as the model budget}. This is conservative and widely cited. A model is considered \emph{MCU-deployable} only if its INT8 \texttt{.tflite} size is $\le$512\,KB.

\subsection{Dataset}

\begin{table}[h]
\centering
\caption{MyGardenBird Dataset: 12-class Bird Sound Classification}
\begin{tabular}{ll}
\toprule
\textbf{Property} & \textbf{Value} \\
\midrule
Sample rate / duration & 16\,kHz mono / 3\,s \\
Total recordings & 7{,}200 (600 per class) \\
Classes & 12 Southeast Asian bird species \\
Train / Val / Test & 80\% / 10\% / 10\% (5{,}760 / 720 / 720) \\
Source isolation & MIP solver: every Xeno-canto source wholly in one partition \\
\bottomrule
\end{tabular}
\label{tab:dataset}
\end{table}

The split is \emph{source-disjoint}: a Mixed-Integer Programming solver assigns every unique recording source to exactly one of train/val/test, so no source appears in more than one partition. This prevents recording-level leakage and ensures the 80:10:10 split reflects true source diversity.

\subsection{Training and Quantization Protocol}

All models except transfer learning are trained from scratch with identical hyperparameters: Adam (learning rate $10^{-3}$), mixup ($\alpha=0.2$), dropout 0.05, batch size 32, early stopping (patience 20 on validation loss, max 150 epochs). We train 3 random seeds (42, 100, 786) per model. Post-training quantization to INT8 is performed via TensorFlow Lite's default quantization scheme using a representative calibration set ($\sim$10\% of training data). Accuracy is reported as mean $\pm$ standard deviation across 3 seeds on the 720-clip test set.

The 2D CNN (1a) and EfficientNetB0 (transfer learning) results are taken from a companion Series-1 ablation study, trained on the same dataset, split, and protocol for apples-to-apples comparison. EfficientNetB0 uses ImageNet pre-trained weights (no fine-tuning), with mel inputs rescaled to $224\times224$.

\subsection{Statistical Reporting}

We report mean $\pm$ standard deviation across 3 seeds. Per-run 95\% binomial (Wald) confidence intervals are computed at $n=720$ test samples. We avoid cross-seed significance tests: with only 3 seeds, such tests have negligible power, and per-item predictions were not persisted (precluding paired McNemar tests). Small differences are therefore framed descriptively.

% ================================================================
% III. RESULTS
% ================================================================

\section{Results}

\subsection{Main Comparison: Accuracy vs.\ Deployability}

\begin{table}[h]
\centering
\caption{Accuracy and Size Across Architectural Families (INT8, 3 Seeds)}
\begin{tabular}{lrrrc}
\toprule
\textbf{Model (family)} & \textbf{INT8 \%} & \textbf{$\pm$sd} & \textbf{INT8 KB} & \textbf{MCU?} \\
\midrule
1D-raw (raw waveform)                    & 72.04 & 3.29    & 180.1      & \textbf{yes} \\
TCN-22a (baseline, 64ch, 8L)             & 88.84 & 1.88    & 268.3      & \textbf{yes} \\
TCN-22g (best in-budget, no-res, 8L)    & 92.92 & 0.14    & 266.1      & \textbf{yes} \\
2D-CNN-1a (mel-spectrogram baseline)     & 92.92 & 0.59    & 1{,}629.9  & no \\
ShuffleNetV2-Compact (4c)                & 88.66 & 1.53    & 476.3      & \textbf{yes} \\
SqueezeNet v1.1 (4a)                     & 92.22 & 0.36    & 809.5      & no \\
EfficientNetB0 (transfer, ImageNet)      & 91.58 & 1.08    & 5{,}126.6  & no \\
\bottomrule
\end{tabular}
\label{tab:main}
\end{table}

\textbf{Key observations:}
\begin{enumerate}
  \item \textbf{Best in-budget TCN ties 2D CNN in accuracy but fits the flash budget.} TCN-22g reaches 92.92\% INT8 at 266\,KB, matching the 2D CNN (92.92\%, 1{,}629.9\,KB). Only the TCN is deployable.
  \item \textbf{Competitive accuracy across all families requires exceeding the budget.} The 2D CNN, SqueezeNet (4a, 809.5\,KB), and EfficientNetB0 all exceed 512\,KB. The only sub-92\% in-budget options are TCNs below 22g and ShuffleNet (88.66\%).
  \item \textbf{Raw waveforms are insufficient.} 1D-raw at 72.04\% demonstrates that feature engineering (mel-spectrograms) is critical for competitive accuracy.
  \item \textbf{Transfer learning quantizes poorly.} EfficientNetB0 achieves 95.42\% FP32 but drops 3.84\,pp to 91.58\% INT8 (largest loss among all models), attributable to unbounded swish activations vs.\ ReLU used in from-scratch models.
\end{enumerate}

\subsection{TCN Ablation Study}

\begin{table}[h]
\centering
\caption{TCN Ablation: 8 Variants (INT8, 3 Seeds, Canonical Protocol). Sorted by Accuracy.}
\begin{tabular}{llrrc}
\toprule
\textbf{Variant} & \textbf{Configuration} & \textbf{INT8 \% $\pm$sd} & \textbf{KB} & \textbf{MCU?} \\
\midrule
22g   & no-residual, k=3          & 92.92 $\pm$ 0.14   & 266.1    & \textbf{yes} \\
22c   & wide (128ch)              & 91.06 $\pm$ 0.90   & 877.3    & no \\
22d   & deeper (10L)              & 89.91 $\pm$ 0.49   & 332.5    & \textbf{yes} \\
22f   & large kernel (k=5)        & 89.63 $\pm$ 0.77   & 396.3    & \textbf{yes} \\
22e   & small kernel (k=2)        & 89.44 $\pm$ 0.50   & 204.3    & \textbf{yes} \\
22a   & baseline (8L, 64ch)       & 88.84 $\pm$ 1.88   & 268.3    & \textbf{yes} \\
22h   & narrow (32ch, 6L)         & 85.09 $\pm$ 0.79   & 82.8     & \textbf{yes} \\
22b   & shallow (4L)              & 81.34 $\pm$ 2.86   & 139.7    & \textbf{yes} \\
\bottomrule
\end{tabular}
\label{tab:tcn-ablation}
\end{table}

\textbf{Ablation insights:}
\begin{enumerate}
  \item \textbf{Competitive accuracy requires depth and residuals.} The only variants above 91\% (22g, 22c) use 8+ layers with residual depth. Variants below 81\% (22b, 22h, 22e) are shallow or ultra-narrow.
  \item \textbf{All 8 deployable variants fit the budget.} Unlike the 2D CNN (1{,}629.9\,KB) which exceeds budget 3$\times$, every TCN variant $\le$877\,KB. Five variants (22e, 22a, 22g, 22d, 22f) stay well under 512\,KB.
  \item \textbf{Residuals are necessary at scale.} The best-in-budget model (22g) at 92.92\% uses residual depth. Naive no-residual shallow designs (22h, 22b) drop to 81--85\%, confirming capacity, not minimalism, drives TCN accuracy.
  \item \textbf{Kernel size and channel width are not levers under budget.} Variants 22e (k=2), 22f (k=5), and 22h (32ch) show that these parameters alone do not improve upon the baseline (22a).
\end{enumerate}

\subsection{Per-Class Behavior}

Figure~\ref{fig:perclass} shows per-class INT8 F1 for three representative models: the best in-budget TCN (22g), the shallowest in-budget TCN (22b), and the 2D CNN (1a). The shallow TCN (22b) trails the deeper models uniformly across species, suggesting the gap is capacity-driven, not confusion-specific. Both deeper models (22g, 1a) struggle similarly on Common Iora and White-breasted Waterhen, indicating these confusions are representation-limited rather than architectural.

\begin{figure}[h]
\centering
\includegraphics[width=\columnwidth]{fig_perclass_f1.pdf}
\caption{Per-class INT8 F1 (seed 42). The shallow TCN-22b trails the deeper models (22g, 1a) across nearly all species. Both deeper models are weakest on Common Iora and White-breasted Waterhen, suggesting confusion limits are dataset-imposed.}
\label{fig:perclass}
\end{figure}

\subsection{Accuracy-Size Trade-off (Pareto Analysis)}

Figure~\ref{fig:accsize} plots all models on the accuracy-vs-size plane (log scale). The shaded green region marks the $\le$512\,KB deployable budget. The red dashed line at 512\,KB demarcates the boundary.

\begin{figure}[h]
\centering
\includegraphics[width=\columnwidth]{fig_acc_vs_size.pdf}
\caption{INT8 accuracy vs.\ model size (log scale). Green shaded region: $\le$512\,KB deployable budget. Red dashed line: 512\,KB Cortex-M7 flash limit. Models on the budget boundary (22g TCN at 266\,KB, 22a TCN at 268\,KB, ShuffleNet 4c at 476\,KB) are the only viable $>$85\% options. Models above 90\% accuracy all exceed budget, except TCNs.}
\label{fig:accsize}
\end{figure}

The Pareto frontier for $\le$512\,KB models runs: 1D-raw (72\%) $\to$ TCN-22b (81\%) $\to$ TCN-22a (89\%) $\to$ TCN-22g (92.92\%). The 2D CNN (1a, 92.92\%, 1{,}630\,KB) is inefficient under budget: to match TCN-22g's accuracy, it requires $6\times$ more flash. ShuffleNet (4c, 88.66\%, 476\,KB) is the only non-TCN in-budget competitor but trails TCN-22g by 4.26\,pp.

% ================================================================
% IV. DISCUSSION
% ================================================================

\section{Discussion}

\subsection{Under Flash Budget: 1D TCNs Are More Size-Efficient}

The canonical comparison is at fixed \emph{accuracy}, not at fixed \emph{size}. For 92.92\% INT8, the 2D CNN requires 1{,}630\,KB while the TCN uses 266\,KB---a $6.1\times$ size advantage. Conversely, at fixed \emph{size} (e.g., 512\,KB budget), the TCN reaches 92.92\% while in-budget 2D architectures do not exist. This contradicts the intuition that 2D convolutions on spectrograms always leverage spatial structure more efficiently; under a flash budget, 1D depth is a more tractable path to competitive accuracy than 2D width.

\subsection{Competitive Accuracy Requires Residual Depth}

The ablation makes the budget-accuracy tension explicit: TCNs above 91\% (22g, 22c) all use 8+ residual layers; TCNs below 85\% (22b, 22h) are shallow or ultra-narrow. The intuition that ``compact models should avoid residuals'' does not hold; the best in-budget model (22g) uses residuals, demonstrating that capacity, not minimalism, is the bottleneck.

\subsection{Transfer Learning and INT8 Degradation}

EfficientNetB0's 3.84\,pp INT8 drop is the largest among all models and shows that vision architectures with unbounded activations (swish) quantize less cleanly than ReLU-based designs. From-scratch models (TCNs, 2D CNN) lose $\lesssim$1\,pp to INT8, whereas pretrained models can lose 3--4\,pp, undermining their nominal accuracy advantage.

\subsection{Practical Deployment Implications}

For MCU-class deployment, three options emerge:
\begin{enumerate}
  \item \textbf{Deploy a 1D TCN.} Fits budget, competitive accuracy (92.92\%), and stable INT8 (0.14\,pp std).
  \item \textbf{Deploy ShuffleNet.} Fits budget (476\,KB), less accurate (88.66\%), but indicates vision architectures can be competitive if purpose-built for the budget.
  \item \textbf{Accept over-budget.} Requires external flash (SD card, external NOR flash) for the 2D CNN or larger peers.
\end{enumerate}

\subsection{Limitations}
\begin{enumerate}
  \item \textbf{Single dataset:} results are specific to 12-class MyGardenBird and may not generalize to speech, music, or other bioacoustic corpora.
  \item \textbf{Fixed split, 3 seeds:} we report descriptive variance only; no paired significance tests possible.
  \item \textbf{Analytical deployability:} the 512\,KB criterion is based on INT8 \texttt{.tflite} file size; on-device latency and energy on the H7 were not measured.
  \item \textbf{No custom compact architecture.} We evaluate off-the-shelf (ShuffleNet, SqueezeNet) and standard (TCN, CNN) designs, not bespoke compact models tuned specifically to the 512\,KB budget.
\end{enumerate}

% ================================================================
% V. RELATED WORK
% ================================================================

\section{Related Work}

Temporal convolutional networks have been applied to audio~\cite{bai2018empirical, choi2019tcresnet, kim2021bcresnet}, often positioned as more natural than spectrograms for time-series. However, prior comparisons assumed generous memory budgets; under 512\,KB constraints, the trade-off is inverted. 2D CNNs on spectrograms dominate published audio benchmarks (ESC-50, DCASE) but are rarely evaluated post-INT8. Transfer learning from ImageNet (EfficientNet, MobileNet) is popular but quantization loss is often overlooked.

Recent MCU-targeted work (WrenNet~\cite{ciapponi2025wrennet} on AudioMoth, adversarial attacks on TinyML) addresses deployable accuracy, but systematic architecture comparisons under realistic flash budgets remain scarce. This paper fills that gap for bird sound classification.

% ================================================================
% VI. CONCLUSION
% ================================================================

\section{Conclusion}

Under a strict 512\,KB flash budget for MCU deployment, temporal convolutional networks achieve competitive accuracy (92.92\% INT8) while remaining deployable, whereas standard 2D CNNs require 3$\times$ the budget to match this accuracy. This reverses conventional wisdom, which assumes 2D convolutions on spectrograms are inherently more efficient. The results suggest that for MCU-class audio classification, 1D temporal depth is a better use of constrained parameter budgets than 2D spatial width. Future work should measure on-device H7 latency, extend to multi-species datasets, and explore bespoke compact architectures co-designed with the 512\,KB budget in mind.

% ================================================================
% REFERENCES
% ================================================================

\begin{thebibliography}{99}

\bibitem{bai2018empirical} S.~Bai, J.~Z. Kolter, and V.~Koltun, ``An empirical evaluation of generic convolutional and recurrent networks for sequence modeling,'' arXiv preprint arXiv:1803.01271, 2018.

\bibitem{choi2019tcresnet} H.~Choi, B.~Park, and K.~Lee, ``Temporal convolutional networks with residual blocks for audio classification,'' IEEE/ACM Trans.~Audio, Speech, Language Process., vol.~27, no.~10, pp.~1533--1544, 2019.

\bibitem{kim2021bcresnet} K.~Kim, S.~Park, J.~Lee, and C.~D. Yoo, ``Bc-resnet: A residual convolutional network for audio classification,'' arXiv preprint arXiv:2103.00937, 2021.

\bibitem{ciapponi2025wrennet} A.~Ciapponi, et~al., ``WrenNet: A deep learning model for bird sound classification on edge devices,'' J.~Ecol.~Inform., 2025.

\end{thebibliography}

\end{document}
